% R. Brown's translations of two letters to L. Breen, with his note and comments (cote 134-2): archive cote 134-2, site page 132 (PDF page 133).
% translation of the letter of 17-19 July 1975. Draft compared: ps1.tex, lines 1538-1570.
% AI-transcribed by Claude Opus 5.5 (claude-opus-5-5, adaptive thinking, effort high), Message Batches API, batch brown_full (msgbatch_01Md3i5yK1PxkpJn5GHfEExE), 2026-11-13; not yet checked by eye.
\scanpage{134-2/132}
\mspage{- 5 -}

Your question 7 seems to indicate that there is discontent in
your mind on the significance of the ``homotopy type'' of $X$ , for
$X$ a topos (for examplethe étale topos of a scheme). Doubtless
you can understand the algebraic topology that one can do on $X$ ;
this includes semi-simplicial sheaves, the stacks of all kinds, and
the relations between these - from this point of view $X$ (with
its very rich structure of topos) virtually disappears so as to
become no more than a pale element of a ``homotopical category''
(or pro-homotopical), deduced from that of the topos by a process
of localisation
- at the first sight, all that still comes from this poor $X$ , are
the $\pi_i$ - and its cohomology groups with constant coefficients -
or at the worst twisted constants. When one digs more into this
definition of ``what is left to this poor $X$'' one falls precisely
on \uline{the} \uline{locally} \uline{constant} \uline{n-stacks}(as an $f : X \to X'$ which is a
homotopy equivalence induces a $(n+1)$-equivalence between the
mentioned locally constant n-stacks on $X$ and on $X'$). - which surely
contains the abelian chain complexes of length $n$ of sheaves with
locally constant cohomology, and the hyper-cohomology of these.
It is thus that one arrives at this triangle of objects which
mutually determine each other :

\diagram{
\begin{tikzcd}[column sep=-3em, every arrow/.append style={<->}]
 & \begin{tabular}{@{}c@{}}
    topos (or topological space\\
    or semi-simplicial complex)\\
    \uline{modulo} n-homotopy
   \end{tabular}
   \arrow[dl]\arrow[dr] & \\
 \begin{tabular}{@{}l@{}}
    n-groupoids\\
    (up to n-equivalence)
   \end{tabular}
   \arrow[rr] &
 & \begin{tabular}{@{}l@{}}
    ``special'' n-categories\\
    (up to n-equivalence).
   \end{tabular}
\end{tikzcd}
}

\mspage{\hand{133}}
